% ============================================================
%  Notes explicatives — Toutes les slides (1 à 21)
%  QMKL-Finance — Mars 2026
% ============================================================
\documentclass[12pt, a4paper]{article}

\usepackage[utf8]{inputenc}
\usepackage[T1]{fontenc}
\usepackage[french]{babel}
\usepackage[margin=2.5cm]{geometry}
\usepackage{titlesec}
\usepackage{xcolor}
\usepackage{mdframed}
\usepackage{booktabs}
\usepackage{amsmath}
\usepackage{amssymb}
\usepackage{hyperref}
\usepackage{enumitem}
\usepackage{parskip}

\definecolor{qblue}{RGB}{10,80,140}
\definecolor{okgreen}{RGB}{0,130,60}
\definecolor{badred}{RGB}{200,40,40}
\definecolor{lightblue}{RGB}{230,240,255}
\definecolor{lightgreen}{RGB}{230,248,235}
\definecolor{lightred}{RGB}{255,235,235}
\definecolor{lightyellow}{RGB}{255,252,230}

\titleformat{\section}
  {\large\bfseries\color{qblue}}
  {\thesection.}
  {8pt}
  {}
  [\vspace{2pt}\hrule\vspace{4pt}]
\titlespacing{\section}{0pt}{18pt}{8pt}

\newmdenv[
  backgroundcolor=lightblue,
  linecolor=qblue, linewidth=2pt,
  leftline=true, rightline=false, topline=false, bottomline=false,
  innerleftmargin=10pt, innerrightmargin=6pt,
  innertopmargin=6pt, innerbottommargin=6pt,
]{retenir}

\newmdenv[
  backgroundcolor=lightred,
  linecolor=badred, linewidth=2pt,
  leftline=true, rightline=false, topline=false, bottomline=false,
  innerleftmargin=10pt, innerrightmargin=6pt,
  innertopmargin=6pt, innerbottommargin=6pt,
]{attention}

\newmdenv[
  backgroundcolor=lightgreen,
  linecolor=okgreen, linewidth=2pt,
  leftline=true, rightline=false, topline=false, bottomline=false,
  innerleftmargin=10pt, innerrightmargin=6pt,
  innertopmargin=6pt, innerbottommargin=6pt,
]{bon}

\newmdenv[
  backgroundcolor=lightyellow,
  linecolor=orange, linewidth=2pt,
  leftline=true, rightline=false, topline=false, bottomline=false,
  innerleftmargin=10pt, innerrightmargin=6pt,
  innertopmargin=6pt, innerbottommargin=6pt,
]{question}

\hypersetup{colorlinks=true, linkcolor=qblue}

% ─────────────────────────────────────────────────────────
\begin{document}

\begin{center}
  {\LARGE\bfseries\color{qblue} Notes explicatives — toutes les slides}\\[6pt]
  {\large Machine Learning Quantique pour la Finance}\\[4pt]
  {\normalsize QMKL sur données réelles · Mars 2026}
\end{center}
\vspace{4pt}\hrule\vspace{10pt}

Ce document accompagne la présentation Beamer (22 slides).
Pour chaque slide : ce qui est montré, pourquoi, et ce qu'il faut retenir
pour répondre aux questions d'un auditoire bancaire ou académique.

\tableofcontents
\newpage

% ═════════════════════════════════════════════════════════════
\section{Slide 1 — Titre}
% ═════════════════════════════════════════════════════════════

\subsection*{Ce que montre la slide}

Page de titre standard : nom du projet, sous-titre, date (Mars 2026).

\subsection*{Ce qu'il faut savoir dire en l'ouvrant}

Cette présentation rend compte d'un projet de recherche appliquée en machine learning quantique
sur des données financières réelles. Le fil conducteur est une question concrète :
\emph{est-ce que les méthodes quantiques apportent un avantage mesurable aujourd'hui ?}
La réponse sera nuancée — honnête sur les limites, optimiste sur les perspectives à 3--5 ans.

Le sigle \textbf{QMKL} signifie \emph{Quantum Multiple Kernel Learning} :
on combine plusieurs "noyaux" quantiques (façons d'encoder les données dans un circuit)
pour obtenir un classifieur robuste.

% ═════════════════════════════════════════════════════════════
\section{Slide 2 — Plan}
% ═════════════════════════════════════════════════════════════

\subsection*{Ce que montre la slide}

Table des matières automatique générée par Beamer.
5 sections : Contexte, Résultats principaux, Comparaison des stratégies,
Notre contribution (QMKL-Boost), Limites, Bilan.

\subsection*{Ce qu'il faut savoir dire}

La présentation suit une logique en entonnoir :
on part de l'hypothèse théorique (pourquoi le quantique ?),
on la teste empiriquement sur 3 datasets financiers,
on analyse les résultats en détail, on présente notre contribution originale (QMKL-Boost),
et on conclut honnêtement sur ce qui est déployable aujourd'hui vs à surveiller.

% ═════════════════════════════════════════════════════════════
\section{Slide 3 — Pourquoi le Machine Learning Quantique en finance ?}
% ═════════════════════════════════════════════════════════════

\subsection*{Ce que montre la slide}

Deux colonnes : l'hypothèse quantique (gauche) et la méthode QMKL (droite).

\subsection*{L'hypothèse quantique}

Un circuit quantique à $Q$ qubits produit un état dans un espace de Hilbert
de dimension $2^Q$. Avec $Q=6$, cela représente un espace de $64$ dimensions ;
avec $Q=20$, plus d'un million de dimensions.
L'idée est que certaines structures dans les données financières
(corrélations entre variables de crédit, de marché, de comportement client)
pourraient être plus facilement capturables dans cet espace de très haute dimension
que par des méthodes classiques à noyau.

\textbf{Nuance importante} : cette propriété n'implique pas automatiquement
un avantage en pratique. Les données financières tabulaires ne sont pas
naturellement "quantiques", et l'espace de Hilbert peut concentrer
les données plutôt que les séparer (voir slide 17 sur les barren plateaus).

\subsection*{La méthode QMKL en 4 étapes}

\begin{enumerate}
  \item Les données client (revenus, historique, etc.) sont encodées dans des qubits
        via des portes de rotation ($R_z$) et d'intrication (CNOT).
  \item Le circuit produit un état quantique $|\psi(x)\rangle$ pour chaque client $x$.
        La similarité entre deux clients est mesurée comme le carré
        de leur "overlap" : $K(x,x') = |\langle\psi(x')|\psi(x)\rangle|^2$.
  \item Un SVM (Support Vector Machine) utilise cette matrice de similarité
        pour apprendre à classer les clients (bon/mauvais payeur, etc.).
  \item On répète cela avec 12 circuits différents (12 kernels)
        et on combine leurs matrices de similarité.
\end{enumerate}

\begin{retenir}
\textbf{Message central :} on ne prétend pas que le quantique est meilleur.
On teste empiriquement si, sur des données réelles d'aujourd'hui,
il apporte un avantage mesurable. C'est une question ouverte.
\end{retenir}

% ═════════════════════════════════════════════════════════════
\section{Slide 4 — Les 3 datasets financiers testés}
% ═════════════════════════════════════════════════════════════

\subsection*{Ce que montre la slide}

Trois blocs décrivent les datasets utilisés, plus la métrique AUC.

\subsection*{Détail des datasets}

\textbf{German Credit} (source : UCI / OpenML, data\_id=31) :\\
1\,000 clients, 20 variables (âge, montant du crédit, durée, type de compte, emploi, etc.),
30\,\% de défauts. C'est le cas d'usage bancaire le plus direct : prédire si un client
va rembourser ou non. Les variables sont hétérogènes (numériques + catégorielles).

\textbf{Bank Marketing} (source : UCI / OpenML, data\_id=1461) :\\
$\sim$45\,000 contacts téléphoniques d'une banque portugaise,
16 variables (durée d'appel, type de contrat, âge, situation professionnelle...).
Seulement 11\,\% de succès : le dataset est très déséquilibré.
Tâche : prédire si un client va souscrire un dépôt à terme.

\textbf{Breast Cancer} (source : sklearn) :\\
569 patients, 30 variables morphologiques (diamètre, texture, concavité des tumeurs).
Très bien séparable (AUC $> 0.99$ pour les classiques). Utilisé comme benchmark
pour mesurer la limite haute atteignable par QMKL dans des conditions favorables.

\subsection*{La métrique AUC}

L'AUC (aire sous la courbe ROC) mesure la capacité du modèle à distinguer
les deux classes, indépendamment du seuil de classification.
\begin{itemize}
  \item AUC = 0.5 : modèle aléatoire, sans information
  \item AUC = 1.0 : classification parfaite
  \item En scoring bancaire, un gain de 1--2 pts d'AUC peut représenter
        des millions d'euros de créances douteuses évitées
\end{itemize}

Tous les résultats sont moyennés sur \textbf{20 répétitions} avec des splits aléatoires
différents pour garantir la robustesse statistique.

% ═════════════════════════════════════════════════════════════
\section{Slide 5 — Du tableau bancaire au circuit quantique : le pipeline}
% ═════════════════════════════════════════════════════════════

\subsection*{Ce que montre la slide}

Le pipeline de prétraitement des données en 3 étapes (gauche),
et la justification de chaque choix (droite).

\subsection*{Étape 1 : Encodage catégoriel}

Les variables textuelles (type d'emploi, type de logement, historique de crédit...)
sont converties en variables binaires 0/1 par \emph{one-hot encoding}
(\texttt{pd.get\_dummies}, \texttt{drop\_first=True} pour éviter la multicolinéarité).
German Credit passe ainsi de 20 variables à 47 colonnes numériques.

\subsection*{Étape 2 : Réduction par PCA}

Un circuit quantique à $Q$ qubits ne peut encoder que $Q$ valeurs simultanément.
Il faut donc réduire la dimension de 47 à $Q$.
La PCA (Analyse en Composantes Principales) sélectionne les $Q$ directions
de variance maximale dans les données.
Avec $Q=6$, on conserve $\approx$85\,\% de la variance totale.

\textbf{Choix de Q :}
\begin{itemize}
  \item Q=4 : plus rapide à simuler, utilisé pour l'analyse principale (German Credit)
  \item Q=6 : meilleur compromis performance/coût pour Bank Marketing et Breast Cancer
  \item Q=8+ : confronte aux barren plateaus (voir slide 17)
\end{itemize}

\subsection*{Étape 3 : Normalisation}

Les angles de rotation des portes quantiques doivent être bornés.
\begin{enumerate}
  \item \textbf{StandardScaler} : centre les données ($\mu=0$, $\sigma=1$)
        pour éviter que les variables à grande variance (revenus, montants)
        dominent l'espace de rotation.
  \item \textbf{MinMaxScaler} : mise à l'échelle vers $[0, 2]$,
        de sorte que l'angle de rotation $\theta = 2\alpha x_i$ explore $[0, 4\pi]$
        — une rotation complète — pour $\alpha \in \{0.5, 4.0\}$.
\end{enumerate}

\begin{retenir}
\textbf{Exemple concret (German Credit, Q=4) :}\\
20 vars $\xrightarrow{\text{one-hot}}$ 47 colonnes
$\xrightarrow{\text{PCA}}$ 4 composantes
$\xrightarrow{\text{scaler}}$ $x_i \in [0,2]$\\
Jeu d'entraînement : 268 clients $\times$ 4 features.
\end{retenir}

% ═════════════════════════════════════════════════════════════
\section{Slide 6 — Comment les données sont encodées dans un circuit quantique}
% ═════════════════════════════════════════════════════════════

\subsection*{Ce que montre la slide}

Le schéma du circuit PauliFeatureMap ZZ pour Q=4 qubits, reps=1,
avec intrication linéaire entre qubits adjacents.

\subsection*{Les 3 blocs du circuit}

\textbf{Bloc 1 — Superposition (portes H, Hadamard) :}\\
Chaque qubit est mis en superposition égale de $|0\rangle$ et $|1\rangle$.
Mathématiquement : $H|0\rangle = \frac{1}{\sqrt{2}}(|0\rangle + |1\rangle)$.
Le circuit explore ainsi simultanément tous les états possibles.

\textbf{Bloc 2 — Encodage individuel ($R_z$) :}\\
Chaque qubit $i$ subit une rotation d'angle $2\alpha x_i$ autour de l'axe Z.
Cela encode la valeur de la $i$-ème feature dans la phase du qubit.

\textbf{Bloc 3 — Encodage des interactions (CNOT + $R_z$) :}\\
Des paires de qubits adjacents $(i, j)$ sont intriquées via des portes CNOT,
puis une rotation $R_z(2\alpha x_i x_j)$ encode le \emph{produit} $x_i x_j$.
Cela capture les corrélations non-linéaires entre features.

\subsection*{La similarité quantique}

$$K(x, x') = |\langle\psi(x')|\psi(x)\rangle|^2$$

C'est le carré du module du produit scalaire entre les états quantiques
produits par les données $x$ et $x'$. Deux clients avec des profils similaires
produisent des états quantiques proches (grande similarité).
C'est l'analogue quantique d'un noyau gaussien RBF.

\begin{question}
\textbf{Question fréquente :} "Pourquoi ZZ et pas juste Z ?"\\
La famille Z encode seulement les effets individuels ($x_i$).
La famille ZZ ajoute les interactions croisées ($x_i x_j$),
ce qui permet de capturer des corrélations non-linéaires entre features.
C'est l'équivalent d'un noyau polynôme de degré 2 vs degré 1.
\end{question}

% ═════════════════════════════════════════════════════════════
\section{Slide 7 — Les 12 kernels quantiques : géométrie de similarité}
% ═════════════════════════════════════════════════════════════

\subsection*{Ce que montre la slide}

6 colonnes (une par famille), chacune avec 3 éléments empilés :
\begin{enumerate}
  \item \textbf{Heatmap 5×5} de la matrice kernel $K(x_i, x_j)$ calculée sur 5 points synthétiques 4D.
  \item \textbf{Badge coloré} indiquant le niveau d'interaction (local / pairwise / 3-corps non-commutatif).
  \item \textbf{Deux mini-barres AUC} (alpha petit et alpha grand) — contexte de performance réelle.
\end{enumerate}

\subsection*{Comment lire une heatmap kernel}

Une heatmap kernel est une représentation de la \textbf{fonction de similarité} dans l'espace quantique.
Chaque case $(i,j)$ représente $K(x_i, x_j) \in [0,1]$ :
\begin{itemize}
  \item \textbf{Case bleue foncée (proche de 1)} : les clients $x_i$ et $x_j$ sont très similaires
        selon ce kernel — ils sont encodés dans des états quantiques proches.
  \item \textbf{Case blanche (proche de 0)} : les clients sont dissemblables —
        leurs états quantiques ont une faible fidélité.
\end{itemize}

Dans nos données synthétiques, les points $x_1$--$x_3$ sont volontairement proches (cluster serré),
et $x_4$--$x_5$ sont lointains. On attend donc un \textbf{bloc bleu en haut à gauche}
(similarité intra-cluster élevée) et des cases pâles dans le coin bas-droite.
Plus ce schéma en blocs est net, plus le kernel discrimine bien les groupes.

\subsection*{La hiérarchie locale / pairwise / 3-corps}

Les 6 familles ne capturent pas le même niveau de géométrie dans les données :

\begin{center}
\begin{tabular}{llll}
  \toprule
  \textbf{Famille} & \textbf{Opérateurs} & \textbf{Niveau} & \textbf{Ce que ça capture} \\
  \midrule
  Z     & $Z$ seul         & Local (1-corps)         & $\cos^2(\alpha(x_i - x_j))$ individuellement \\
  ZZ    & $Z$, $ZZ$        & Pairwise (2-corps)      & + termes croisés $x_k \cdot x_l$ \\
  XZ    & $X$, $Z$         & Pairwise (2-corps)      & Rotations dans le plan XZ \\
  YXX   & $Y$, $XX$        & 3-corps non-comm.       & Corrélations quadratiques $x^2$ \\
  YZX   & $Y$, $ZX$        & 3-corps non-comm.       & Interactions non-commutatives croisées \\
  Pauli & $Z$, $ZZ$        & Pairwise (2-corps)      & ZZ à faible bande passante ($\alpha \times 0.6$) \\
  \bottomrule
\end{tabular}
\end{center}

\textbf{Intuition :} un kernel local (Z) produit des heatmaps avec une transition progressive
— il "voit" chaque feature indépendamment. Un kernel 3-corps (YXX, YZX) peut créer des
discontinuités plus nettes, au prix d'une plus grande sensibilité au bruit.

\subsection*{Le paramètre $\alpha$ (largeur de bande)}

$\alpha$ est l'hyperparamètre d'échelle du kernel quantique.
Il joue exactement le même rôle que $\gamma$ dans un kernel RBF classique
$K_{\text{RBF}}(x, x') = \exp(-\gamma \|x-x'\|^2)$ :
\begin{itemize}
  \item $\alpha$ petit (0.5--1.0) : bande large, le kernel varie lentement.
        Deux clients peuvent être distants et rester similaires.
        Les mini-barres AUC montrent les valeurs pour chaque $\alpha$.
  \item $\alpha$ grand (2.5--4.0) : bande étroite, le kernel est très sélectif.
        Seuls les clients très proches ont une similarité élevée.
\end{itemize}

\subsection*{Lien avec l'analyse Shapley (Slide 14)}

La valeur de Shapley de chaque kernel mesure sa contribution marginale à la combinaison QMKL.
Les familles YXX et ZZ ont les Shapley les plus élevés ($>+0.03$) parce qu'elles apportent
une \textbf{information géométrique complémentaire} : ZZ capture les corrélations paires
absentes de Z ; YXX capture les non-linéarités quadratiques absentes de ZZ.

\begin{retenir}
\textbf{Pourquoi 12 kernels et pas 1 ?}\\
Aucun kernel seul n'est optimal sur tous les datasets.
La hiérarchie local / pairwise / 3-corps garantit que les 6 familles explorent
des géométries différentes dans l'espace de Hilbert quantique.
QMKL-Boost (QUBO) sélectionne les 2 familles les plus complémentaires.
\end{retenir}

% ═════════════════════════════════════════════════════════════
\section{Slide 8 — Comment fonctionne chaque stratégie QMKL ?}
% ═════════════════════════════════════════════════════════════

\subsection*{Ce que montre la slide}

Cinq blocs décrivent les cinq stratégies de combinaison testées.

\subsection*{Explication de chaque stratégie}

\textbf{Average — Moyenne uniforme :}\\
$K_{\text{mkl}} = \frac{1}{M}\sum_{m=1}^{M} K_m$.
Chaque kernel contribue à parts égales, aucun paramètre à régler.
Intuitif mais aveugle à la qualité individuelle de chaque kernel.
Souffre quand $N$ est faible et $Q$ petit.

\textbf{Centered Alignment :}\\
Maximise analytiquement $\langle K_{\text{mkl}}, yy^\top \rangle_F$,
la corrélation de Frobenius entre la matrice kernel combinée et les labels.
Solution en forme fermée, sans itération ni hyperparamètre.
S'adapte à la structure du dataset.

\textbf{Bayesian Optimization :}\\
Poids $w_m \in [0,1]$ optimisés par cross-validation en boîte noire (processus gaussien).
200 évaluations = 200$\times$ plus lent que Centered Alignment.
Gain marginal ne justifiant pas le coût.

\textbf{QMKL-Boost (notre contribution — QUBO) :}\\
Sélection binaire $w_m \in \{0,1\}$ via un problème QUBO (Quadratic Unconstrained Binary Optimization).
Sélectionne les 2 kernels les plus complémentaires parmi 12.
Formulation nativement compatible avec les annealers quantiques.

\textbf{VQKL — Variationnel :}\\
Poids appris par descente de gradient.
Avec $N=40$--200, sur-apprend sans converger vers une solution stable.

\begin{retenir}
Centered Alignment est la stratégie la plus efficace sur German Credit (Q=4, N=40).
La complexité algorithmique ne garantit pas de meilleures performances
avec de petits datasets.
\end{retenir}

% ═════════════════════════════════════════════════════════════
\section{Slide 9 — Résultat 1 : QMKL vs méthodes classiques}
% ═════════════════════════════════════════════════════════════

\subsection*{Ce que montre la slide}

Graphique en barres groupées : QMKL-Centered vs RBF-SVM vs Forêt aléatoire, pour chaque dataset.

\subsection*{Lecture des résultats}

\begin{center}
\begin{tabular}{lccc}
  \toprule
  \textbf{Dataset} & \textbf{QMKL-Centered} & \textbf{RBF-SVM} & \textbf{Forêt aléatoire} \\
  \midrule
  German Credit  & 0.743 & 0.800 & 0.845 \\
  Bank Marketing & 0.861 & 0.875 & 0.890 \\
  Breast Cancer  & 0.993 & 0.997 & 0.998 \\
  \bottomrule
\end{tabular}
\end{center}

\textbf{German Credit :} écart de $-6$ pts vs RBF-SVM. Significatif en scoring bancaire.\\
\textbf{Bank Marketing :} quasi-parité ($-1.4$ pts). Dataset plus favorable à QMKL.\\
\textbf{Breast Cancer :} quasi-parité ($-0.4$ pts). Données bien séparables.

\begin{attention}
\textbf{Ne pas généraliser :} le bon résultat de Breast Cancer s'explique par la séparabilité naturelle
des données, pas par un avantage structurel du quantique. Sur les vrais cas financiers difficiles,
l'écart est systématiquement en faveur des classiques.
\end{attention}

% ═════════════════════════════════════════════════════════════
\section{Slide 10 — Résultat 2 : Classement détaillé sur German Credit}
% ═════════════════════════════════════════════════════════════

\subsection*{Ce que montre la slide}

Graphique horizontal avec toutes les méthodes classées par AUC décroissante
(German Credit, Q=4, N=40).

\subsection*{Analyse du classement}

\textbf{Top 3 — méthodes classiques :}
\begin{enumerate}
  \item Forêt aléatoire : AUC = 0.845
  \item RBF-SVM : AUC = 0.800
  \item Régression logistique : AUC = 0.783
\end{enumerate}

\textbf{Meilleur QMKL — 4\textsuperscript{e} place :}\\
Centered Alignment (0.743), à $-6$ pts du leader.

\textbf{Reste du classement QMKL :}
QKRR (0.714), VQKL (0.688), BO (0.672), Average (0.635), Single-Best (0.575).

Les méthodes plus complexes (QKRR, BO) ne rattrapent pas le leader QMKL (Centered).
Single-Best confirme qu'un kernel seul est insuffisant.

\begin{retenir}
Pour le risque crédit en production aujourd'hui : RBF-SVM ou Forêt aléatoire.
Aucune méthode QMKL ne dépasse le seuil classique sur ce dataset.
\end{retenir}

% ═════════════════════════════════════════════════════════════
\section{Slide 11 — Résultat 3 : Plus de données n'efface pas l'écart}
% ═════════════════════════════════════════════════════════════

\subsection*{Ce que montre la slide}

Courbes d'apprentissage : AUC en fonction de $N$ (40 à 200 clients),
pour QMKL-Centered, RBF-SVM, et Forêt aléatoire sur German Credit, Q=4.

\subsection*{Hypothèse testée}

Un déficit de données d'entraînement pourrait-il expliquer l'écart QMKL / classiques ?
Si oui, ajouter des données résoudrait le problème.

\subsection*{Résultat}

Non. De $N=40$ à $N=200$, les trois méthodes progressent en parallèle à la même vitesse.
L'écart de $\approx 7$ pts reste constant sur toute la plage.
Les pentes de progression sont similaires ($\approx 0.16$ pt d'AUC par 100 clients supplémentaires).

\subsection*{Interprétation}

Il s'agit d'une \textbf{limite structurelle} des circuits PauliFeatureMap sur données tabulaires financières,
pas d'un manque de données. La forme des kernels quantiques (PauliFeatureMap ZZ)
ne capture pas les mêmes structures que les kernels RBF classiques sur ce type de données,
indépendamment de $N$.

\begin{attention}
\textbf{Implication R\&D :} collecter plus de données ne résoudra pas le problème.
Des architectures de circuits différentes ou des méthodes d'encodage adaptatives
seraient nécessaires pour combler cet écart.
\end{attention}

% ═════════════════════════════════════════════════════════════
\section{Slide 12 — Les 5 stratégies QMKL face à face}
% ═════════════════════════════════════════════════════════════

\subsection*{Ce que montre la slide}

Trois graphiques (un par dataset) avec l'AUC de chaque stratégie.

\subsection*{Analyse par dataset}

\textbf{German Credit (Q=4, N=40) :}
Centered domine (0.743). Average s'effondre (0.635) car les kernels peu informatifs
bruitent la combinaison. Single-Best frôle l'aléatoire (0.575).

\textbf{Bank Marketing (Q=6, N=100) :}
Centered reste le meilleur (0.861). L'écart avec BO (0.793) se réduit
car $N$ plus grand aide BO à converger. Average reste médiocre (0.612).

\textbf{Breast Cancer (Q=6, N=100) :}
Toutes les stratégies dépassent 0.92. Les données sont très séparables —
même Average performe bien (0.958). Ce cas ne discrimine pas les stratégies.

\begin{retenir}
Centered Alignment est la stratégie la plus robuste : jamais la pire,
souvent la meilleure, sans aucun hyperparamètre à régler.
\end{retenir}

% ═════════════════════════════════════════════════════════════
\section{Slide 13 — Quelle stratégie QMKL choisir ?}
% ═════════════════════════════════════════════════════════════

\subsection*{Ce que montre la slide}

Tableau AUC / coût de calcul + 3 enseignements clés (German Credit, Q=4, N=40).

\subsection*{Analyse du tableau}

\begin{center}
\begin{tabular}{lcc}
  \toprule
  \textbf{Stratégie} & \textbf{AUC} & \textbf{Coût calcul} \\
  \midrule
  \textbf{Centered Alignment} & \textbf{0.743} & Très faible \\
  QKRR              & 0.714 & Faible \\
  VQKL              & 0.688 & Moyen \\
  Bay. Optim.       & 0.672 & 200$\times$ plus cher \\
  Average           & 0.635 & Nul \\
  \bottomrule
\end{tabular}
\end{center}

\textbf{Pourquoi Average est-il le pire ici ?}
Avec Q=4 et N=40, plusieurs kernels sont peu informatifs.
La moyenne uniforme les dilue avec les bons kernels, dégradant le signal.

\textbf{Pourquoi BO n'est-il pas recommandé ?}
200 évaluations de cross-validation pour un gain de 4 pts vs Average
(et encore $-$7 pts vs Centered) : coût prohibitif en production.

\textbf{Pourquoi VQKL ne converge pas ?}
Le gradient descent sur les poids oscille sans trouver de minimum stable
avec $N=40$. Sur-apprentissage garanti.

\begin{retenir}
\textbf{Recommandation :} Centered Alignment sur risque crédit (Q=4).
Rapide, stable, sans tuning, adaptatif à la structure du dataset.
\end{retenir}

% ═════════════════════════════════════════════════════════════
\section{Slide 14 — Quels kernels apportent le plus ? (Analyse Shapley)}
% ═════════════════════════════════════════════════════════════

\subsection*{Ce que montre la slide}

Deux graphiques : valeurs de Shapley par kernel (gauche)
et poids Centered Alignment pour les kernels positifs (droite), sur German Credit.

\subsection*{Qu'est-ce que la valeur de Shapley ?}

Concept de théorie des jeux coopératifs : la valeur de Shapley $\phi_m$ d'un kernel $m$
est sa contribution marginale \emph{moyenne} sur toutes les $2^{12}-1 = 4095$ coalitions possibles.
Un kernel avec $\phi_m > 0$ améliore la combinaison ; avec $\phi_m < 0$, il la dégrade.

\subsection*{Résultats sur German Credit}

\textbf{Top contributeurs :} Z $\alpha$=1.0 ($\phi = +0.039$, poids CA = 0.305),
YZX $\alpha$=0.6 ($\phi = +0.031$, poids = 0.248), YXX $\alpha$=0.6 ($\phi = +0.018$).

\textbf{Kernels nuisibles :} Pauli $\alpha$=2.5 ($\phi = -0.017$), XZ $\alpha$=2.5 ($\phi = -0.006$).
Les retirer améliorerait QMKL — ce que fait QMKL-Boost.

\textbf{Pourquoi Z $\alpha$=1.0 domine-t-il ?}
Il est le plus \emph{différent} des autres structurellement : il encode uniquement les effets
individuels (pas d'interactions croisées) avec une bande large.
Il est donc complémentaire aux kernels ZZ et YZX.

\begin{bon}
\textbf{Corrélation diversité $\leftrightarrow$ gain :} $r = 0.74$
entre la dissimilarité d'un kernel aux autres et sa valeur de Shapley.
Plus un kernel est différent, plus il contribue.
L'avantage QMKL vient de la \textbf{diversité}, pas du nombre de kernels.
\end{bon}

% ═════════════════════════════════════════════════════════════
\section{Slide 15 — QMKL-Boost : sélection intelligente des kernels}
% ═════════════════════════════════════════════════════════════

\subsection*{Ce que montre la slide}

Description de QMKL-Boost + graphique comparant toutes les méthodes sur 2 datasets.

\subsection*{Le problème d'Average}

Moyenner 12 kernels uniformément dilue les kernels informatifs avec les kernels
redondants ou nuisibles (confirmé par l'analyse Shapley).
Le signal utile est noyé dans le bruit.

\subsection*{La solution QUBO}

$$\min_{w \in \{0,1\}^M}
  \underbrace{-\sum_m w_m s_m}_{\text{performance individuelle}}
  + \lambda \underbrace{\sum_{m \neq m'} w_m w_{m'} A_{mm'}}_{\text{pénalité redondance}}$$

\begin{itemize}
  \item $s_m$ : AUC individuelle du kernel $m$ sur le jeu d'entraînement
  \item $A_{mm'}$ : corrélation entre les matrices kernel $m$ et $m'$
  \item $\lambda$ : compromis performance / diversité (à calibrer par sweep)
\end{itemize}

\subsection*{Résultats}

\begin{center}
\begin{tabular}{lcc}
  \toprule
  \textbf{Méthode} & \textbf{German Credit} & \textbf{Bank Marketing} \\
  \midrule
  Average (12 kernels)     & 0.782 & 0.691 \\
  Centered Alignment        & 0.799 & 0.730 \\
  \textbf{QUBO (2 kernels)} & \textbf{0.823} & \textbf{0.767} \\
  \midrule
  Gain vs Average           & \textbf{+4.1 pts} & \textbf{+7.6 pts} \\
  \bottomrule
\end{tabular}
\end{center}

Kernels sélectionnés : \texttt{ZZ $\alpha$=2.5} + \texttt{XZ $\alpha$=0.5}
(deux familles structurellement différentes).

\begin{bon}
\textbf{Connexion QPU :} QUBO est le format natif des annealers quantiques
(D-Wave, Pasqal, IBM). Notre formulation est directement transposable sur hardware
quantique quand il sera suffisamment stable.
\end{bon}

% ═════════════════════════════════════════════════════════════
\section{Slide 16 — La matrice Q — comment QUBO choisit les kernels}
% ═════════════════════════════════════════════════════════════

\subsection*{Ce que montre la slide}

Deux éléments côte à côte :
\begin{enumerate}
  \item \textbf{Heatmap 12×12 de la matrice Q QUBO} — couleurs RdBu centrées à 0.
  \item \textbf{Tableau de décision conceptuel} (4 cases : AUC haute/basse × diversifiés/redondants).
\end{enumerate}

\subsection*{Lecture de la matrice Q (diagonale vs hors-diagonale)}

La matrice Q est la représentation matricielle du problème QUBO :
$$\min_{w \in \{0,1\}^M} w^\top Q\, w$$

\begin{itemize}
  \item \textbf{Diagonale $Q_{mm} = -s_m < 0$} (cases bleues) :\\
        Un terme négatif \emph{encourage} le solveur à mettre $w_m = 1$.
        Plus $s_m$ est élevé (bonne AUC individuelle), plus la case est bleu foncé
        — le kernel est attractif.

  \item \textbf{Hors-diagonale $Q_{mm'} = \lambda A_{mm'} > 0$} (cases rouges) :\\
        Un terme positif \emph{pénalise} la co-sélection de deux kernels redondants.
        $A_{mm'}$ est l'alignement de Frobenius entre les matrices kernel $m$ et $m'$.
        Même famille = $A \approx 0.80$ (très rouge) ;
        familles croisées différentes = $A \approx 0.20$ (presque blanc).
\end{itemize}

Le solveur QUBO équilibre ces deux forces : sélectionner les meilleurs kernels
(maximiser la somme des $s_m$) tout en évitant les paires redondantes
(minimiser la somme des $A_{mm'}$).

\subsection*{Tableau de décision}

Le tableau 2×2 synthétise la logique de sélection :

\begin{center}
\begin{tabular}{lcc}
  \toprule
  & \textbf{AUC haute} & \textbf{AUC basse} \\
  \midrule
  \textbf{Diversifiés} ($A < 0.4$)  & SÉLECTIONNER \checkmark & Peut-être (contexte) \\
  \textbf{Redondants} ($A > 0.7$)   & Choisir 1/2    & IGNORER $\times$ \\
  \bottomrule
\end{tabular}
\end{center}

\textbf{Résultat à $\lambda^*=1$ :} ZZ $\alpha$=4.0 + XZ $\alpha$=0.5.
Ces deux kernels sont dans la case "SÉLECTIONNER" : bonne AUC individuelle ($s > 0.72$)
et faible redondance mutuelle ($A \approx 0.20$).

\subsection*{Calibration de $\lambda$}

Un sweep sur 8 valeurs de $\lambda$ (8 entraînements SVM) suffit.
La courbe AUC vs $\lambda$ est unimodale : maximum à $\lambda=1$, décroissante au-delà.
La figure \texttt{pres\_F10\_lambda\_sweep.png} (conservée dans \texttt{results/presentation/})
détaille cette courbe.

\begin{retenir}
\textbf{Gain de production :} 2 kernels = 6$\times$ moins de calcul de matrice kernel
que 12 kernels. Latence de scoring réduite d'un facteur 6 à l'inférence.
\end{retenir}

% ═════════════════════════════════════════════════════════════
\section{Slide 16b — Quand le quantique devient-il indispensable pour QUBO ?}
% ═════════════════════════════════════════════════════════════

\subsection*{Ce que montre la slide}

Graphique en log-scale comparant trois courbes de complexité en fonction de $M$ (nombre de kernels) :
\begin{itemize}
  \item Brute force classique $O(2^M)$ — rouge tiretté
  \item Heuristique classique $O(M^3)$ — orange pointillé
  \item Recuit quantique $O(M^{1.5})$ empirique — bleu plein
\end{itemize}
Trois zones colorées : classique (M < 25), transition (25–35), quantique (M > 35).

\subsection*{Pourquoi ce slide est important}

Sans ce slide, l'audience peut se demander : "pourquoi QUBO ? un simple tri suffit."
Ce graphique montre que la sélection optimale de kernels est un \textbf{problème NP-difficile} :
l'espace de recherche est $2^M$ sous-ensembles, qui explose de manière exponentielle.

\begin{center}
\begin{tabular}{rll}
  \toprule
  \textbf{M} & \textbf{Sous-ensembles $2^M$} & \textbf{Remarque} \\
  \midrule
  12  & 4\,096           & Notre projet — classique parfaitement viable \\
  25  & 33 millions      & Limite classique (quelques minutes CPU) \\
  35  & 34 milliards     & Zone de transition \\
  50  & $\approx 10^{15}$ & CPU : des millénaires ; quantique : quelques secondes \\
  \bottomrule
\end{tabular}
\end{center}

\subsection*{Honnêteté académique}

Le recuit quantique (D-Wave, Pasqal) est présenté comme \textbf{heuristique empiriquement polynomial}.
Il n'existe pas de preuve théorique que le recuit quantique soit plus rapide que le recuit simulé
classique pour QUBO de petite taille. La complexité $O(M^{1.5})$ est une observation empirique
sur des benchmarks publics, pas une borne garantie.

\begin{retenir}
\textbf{Ce n'est pas la simulation quantique — c'est l'optimisation combinatoire.}\\
Ce slide corrige une confusion fréquente : le quantique n'est pas utile ici pour
calculer les matrices kernel (c'est classique), mais pour \emph{optimiser la sélection}
quand $M$ devient grand ($M \geq 25$).
\end{retenir}

\subsection*{Questions fréquentes sur ce slide}

\begin{question}
\textbf{Q : D-Wave est-il disponible commercialement ?}\\
A : Oui, D-Wave propose un accès cloud (Leap). Pasqal est une startup française avec
accès via API. IBM propose également des optimiseurs QUBO sur ses QPU via QAOA.
Mais pour $M=12$, le classique est parfaitement suffisant — le slide montre à partir
de \emph{quand} le quantique devient nécessaire, pas que nous en avons besoin aujourd'hui.
\end{question}

\begin{question}
\textbf{Q : Pourquoi ne pas utiliser une heuristique classique (recuit simulé, algorithmes génétiques) ?}\\
A : Pour $M \leq 35$, c'est exactement ce que nous faisons — la courbe $O(M^3)$ représente
ces heuristiques. Elles sont très efficaces. La valeur du recuit quantique n'apparaît qu'au-delà
du break-even $M \approx 25$--35. Notre projet à $M=12$ n'en a pas besoin, mais la
\emph{formulation QUBO} reste directement transposable sans modification.
\end{question}

% ═════════════════════════════════════════════════════════════
\section{Slide 18 — La limite fondamentale : barren plateaus}
% ═════════════════════════════════════════════════════════════

\subsection*{Ce que montre la slide}

Deux graphiques : variance du kernel vs nombre de qubits $Q$,
et heatmap de concentration par kernel et par $Q$.

\subsection*{Définition des barren plateaus}

À mesure que $Q$ augmente, les valeurs du kernel quantique
$K(x, x') = |\langle\psi(x')|\psi(x)\rangle|^2$ se concentrent exponentiellement
vers une constante (loi de la concentration de la mesure).
Les données deviennent \textbf{indiscernables} pour le circuit :
deux clients très différents ont le même score de similarité quantique.

\subsection*{Mesure de l'effet}

\begin{center}
\begin{tabular}{lcc}
  \toprule
  \textbf{Q} & \textbf{Variance kernels globaux (ZZ)} & \textbf{Variance kernels locaux (Z)} \\
  \midrule
  2 & 1.000 & 1.000 \\
  4 & 0.650 & 0.760 \\
  6 & 0.450 & 0.590 \\
  8 & 0.400 & 0.500 \\
  \bottomrule
\end{tabular}
\end{center}

Les kernels locaux (Z seul, sans interactions croisées) perdent 50\,\% de variance
contre 60\,\% pour les kernels globaux (ZZ) — ils sont 40\,\% plus robustes.

\subsection*{Conséquences pratiques}

\begin{attention}
\textbf{Pour la banque :} augmenter $Q$ ne garantit pas de meilleures performances.
C'est un problème ouvert en recherche. Pistes de mitigation connues :
utiliser des kernels locaux (Z), réduire la profondeur du circuit (reps=1),
ou apprendre l'encodage (feature maps adaptatives).
\end{attention}

% ═════════════════════════════════════════════════════════════
\section{Slide 19 — Test sur IBM Torino : kernels réels vs simulation}
% ═════════════════════════════════════════════════════════════

\subsection*{Ce que montre la slide}

Comparaison AUC simulateur statevector vs hardware IBM Torino,
pour 3 kernels sur German Credit (N=30, Q=4).

\subsection*{Setup expérimental}

\begin{itemize}
  \item Backend : \textbf{IBM Torino} (processeur Heron r2, 133 qubits supraconducteurs)
  \item Shots : 1\,024 mesures par circuit (estimation de la probabilité par Monte Carlo)
  \item N=30 clients (petit $N$ pour limiter le coût QPU)
  \item 3 kernels testés : ZZ($\alpha$=1), ZZ($\alpha$=4), XZ($\alpha$=0.5)
\end{itemize}

\subsection*{Observations}

\textbf{ZZ($\alpha$=1)} : hardware \emph{meilleur} que la simulation (+6.5 pts AUC).
Le bruit du QPU agit comme une régularisation implicite sur ce kernel.
C'est un effet ambigu : une amélioration non contrôlée n'est pas fiable.

\textbf{ZZ($\alpha$=4)} : AUC $< 0.50$ sur hardware. Le kernel est trop concentré
(fort barren plateau) ; le bruit l'effondre complètement.

\textbf{XZ($\alpha$=0.5)} : légère dégradation sur hardware ($-$2.5 pts).
Comportement différent des deux cas précédents.

\begin{attention}
\textbf{Verdict :} résultats incohérents d'un kernel à l'autre.
Le bruit QPU améliore certains kernels et en détruit d'autres —
comportement imprévisible, non contrôlable en production.
\end{attention}

% ═════════════════════════════════════════════════════════════
\section{Slide 20 — IBM Torino : robustesse des stratégies MKL face au bruit}
% ═════════════════════════════════════════════════════════════

\subsection*{Ce que montre la slide}

Tableau comparant l'AUC simulateur vs hardware pour 4 stratégies QMKL + RBF-SVM.

\subsection*{Analyse}

\begin{center}
\begin{tabular}{lcc}
  \toprule
  \textbf{Méthode} & \textbf{AUC Simulateur} & \textbf{AUC Hardware} \\
  \midrule
  \textbf{Centered Alignment} & 0.735 & \textbf{0.735} \\
  Bayesian Optim.             & 0.580 & 0.670 \\
  Average                     & 0.575 & 0.475 \\
  Projection                  & 0.435 & 0.225 \\
  \midrule
  RBF-SVM (classique)         & \multicolumn{2}{c}{0.775} \\
  \bottomrule
\end{tabular}
\end{center}

\textbf{Centered Alignment} est la seule stratégie parfaitement stable :
AUC identique en simulation et sur hardware (0.735).
Sa nature analytique (poids calculés sur la matrice kernel, pas sur les données brutes)
la rend robuste aux erreurs de mesure.

\textbf{Average} perd 10 pts sur hardware : la moyenne est sensible aux erreurs
sur chaque kernel, et les erreurs s'accumulent.

\textbf{Projection} s'effondre de 0.435 à 0.225 : cette stratégie amplifie le bruit.

\begin{retenir}
Même sur QPU réel, QMKL-Centered (0.735) reste $-$4 pts sous RBF-SVM classique (0.775).
Centered Alignment est la stratégie recommandée si on doit tester sur hardware.
\end{retenir}

% ═════════════════════════════════════════════════════════════
\section{Slide 21 — Ce que ça veut dire pour la banque}
% ═════════════════════════════════════════════════════════════

\subsection*{Ce que montre la slide}

Deux colonnes : verdict aujourd'hui (ne pas déployer) vs à surveiller (horizon 3--5 ans).

\subsection*{Aujourd'hui — Ne pas déployer}

\begin{itemize}
  \item \textbf{Performance insuffisante :} $-7$ pts AUC sur le risque crédit.
        RBF-SVM et Forêt aléatoire sont supérieurs, moins coûteux et plus rapides.
  \item \textbf{Hardware trop bruité :} variance $\pm 0.20$ entre runs sur QPU réel.
        Inacceptable pour des décisions de crédit en production.
  \item \textbf{Limite structurelle :} plus de données n'efface pas l'écart.
        Le problème n'est pas statistique, il est architectural.
\end{itemize}

\subsection*{À surveiller — Horizon 2028--2030}

\begin{itemize}
  \item \textbf{QMKL-Boost (+4 pts via QUBO) :} notre approche de sélection de kernels
        dépasse Centered Alignment. Quand le hardware sera stable,
        ce gain sera exploitable directement sur QPU.
  \item \textbf{Interprétabilité réglementaire (Bâle IV, RGPD) :}
        QUBO sélectionne 2 kernels sur 12, ce qui produit un modèle
        plus explicable qu'une combinaison de 12 noyaux opaques.
  \item \textbf{Vitesse d'inférence :} 2 kernels = 6$\times$ moins de calcul
        à l'inférence, avantage potentiel pour le scoring temps réel.
\end{itemize}

\begin{retenir}
\textbf{Position recommandée :} R\&D exploratoire à budget contenu.
Réévaluation complète en 2028--2030 quand la correction d'erreur QPU sera mature.
Ne pas investir dans du hardware QPU pour la classification de crédit à court terme.
\end{retenir}

% ═════════════════════════════════════════════════════════════
\section{Slide 22 — Merci / Questions}
% ═════════════════════════════════════════════════════════════

\subsection*{Ce que montre la slide}

Slide de clôture avec le message de synthèse : \\
\emph{"QMKL-Boost +4 pts AUC avec 2 kernels sur 12 — l'avantage quantique est dans la sélection, pas dans le nombre."}

\subsection*{Questions fréquentes à anticiper}

\begin{question}
\textbf{Q : Pourquoi pas tester sur plus de qubits ?}\\
A : Au-delà de Q=8, les barren plateaus dégradent significativement les kernels.
La variance utile perd 60\,\% entre Q=2 et Q=8. Augmenter Q sans corriger ce phénomène
est contre-productif. Des recherches actives visent à le résoudre (kernels locaux, encodages adaptatifs).
\end{question}

\begin{question}
\textbf{Q : Ces résultats sont-ils publiables / reproductibles ?}\\
A : Oui. 9 notebooks exécutés, 20 runs par expérience, splits aléatoires documentés.
Les données sont publiques (OpenML). Les kernels sont implémentés via Qiskit 1.x
et qiskit-machine-learning.
\end{question}

\begin{question}
\textbf{Q : Quand sera-t-il raisonnable de déployer du QMKL en production ?}\\
A : Quand la correction d'erreur quantique (QEC) sera opérationnelle à grande échelle
— estimée vers 2028--2030 selon IBM et Google. D'ici là, les algorithmes variationnels
et le QUBO classique restent les options les plus réalistes.
\end{question}

\begin{question}
\textbf{Q : Pourquoi Centered Alignment est-il aussi robuste ?}\\
A : Parce que c'est une solution analytique en forme fermée — elle ne dépend pas
d'un processus d'optimisation itératif qui pourrait sur-apprendre ou diverger.
Les poids sont calculés directement à partir de la matrice kernel
et des labels d'entraînement, sans tuning.
\end{question}

% ═════════════════════════════════════════════════════════════
\section*{Annexe — Questions classiques à anticiper}
\addcontentsline{toc}{section}{Annexe — Questions classiques à anticiper}
% ═════════════════════════════════════════════════════════════

Cette annexe regroupe les questions les plus fréquentes dans ce type de présentation,
organisées par thème. Pour chaque question : une réponse directe, le chiffre à citer,
et la slide de référence.

% ─────────────────────────────────────────────────────────
\subsection*{A. Choix du nombre de qubits Q}
% ─────────────────────────────────────────────────────────

\begin{question}
\textbf{Q : Pourquoi choisir Q=4 et pas Q=10 ou Q=20 pour exploiter davantage le quantique ?}\\[4pt]
\textbf{Réponse :} Deux raisons opposées contraignent $Q$ :
\begin{enumerate}
  \item \textbf{Barren plateaus (trop grand Q) :} au-delà de Q=6--8, la variance des kernels
        s'effondre exponentiellement. La matrice kernel devient presque constante —
        les données sont indiscernables. Aucun gain AUC ne compense cette perte d'information.
        On mesure $-60\,\%$ de variance entre Q=2 et Q=8 sur ZZ.
  \item \textbf{Mismatch de dimension (trop petit Q) :} il faut au moins autant de qubits
        que de features utiles. Avec Q=4, la PCA conserve $\approx 75\,\%$ de la variance
        de German Credit. Descendre à Q=2 ferait perdre trop d'information.
\end{enumerate}
Q=4--6 est donc le \emph{sweet spot} empirique pour données tabulaires financières (2026).
\hfill\textit{→ Slide 17}
\end{question}

\begin{question}
\textbf{Q : Est-ce qu'augmenter Q améliorerait les performances si le hardware était parfait ?}\\[4pt]
\textbf{Réponse :} Pas nécessairement. Les barren plateaus sont un phénomène théorique
lié à la structure des circuits de type PauliFeatureMap, indépendant du bruit hardware.
Même sur simulateur parfait (statevector), la variance des kernels chute avec Q.
Des architectures alternatives (circuits locaux, encodage adaptatif) pourraient mitiger ce problème
mais ne sont pas encore comparables aux classiques sur données tabulaires.
\hfill\textit{→ Slide 17}
\end{question}

\begin{question}
\textbf{Q : Pourquoi ne pas utiliser plus de répétitions (\texttt{reps}) du circuit ?}\\[4pt]
\textbf{Réponse :} Chaque répétition ajoute une couche d'intrication.
Plus de répétitions = circuits plus profonds = plus sensibles au bruit QPU,
et variance encore plus concentrée (les barren plateaus empirent).
Sur simulateur, 2 reps vs 1 rep ne change pas significativement l'AUC
mais double la profondeur de circuit et le coût de calcul.
\hfill\textit{→ Slides 6, 17}
\end{question}

% ─────────────────────────────────────────────────────────
\subsection*{B. Choix des familles de kernels et du paramètre $\alpha$}
% ─────────────────────────────────────────────────────────

\begin{question}
\textbf{Q : Comment avez-vous choisi les 6 familles de kernels ?}\\[4pt]
\textbf{Réponse :} Les familles ont été sélectionnées pour couvrir l'espace des opérateurs de Pauli
(Z, X, Y et leurs produits tensoriels) avec une diversité structurelle maximale.
Z capture uniquement les effets individuels ; ZZ ajoute les interactions paires ;
YXX/YZX ajoutent des interactions à 3 corps et des rotations non-commutatives.
L'objectif est que les 12 kernels soient complémentaires — la corrélation $r=0.74$
entre dissimilarité structurelle et valeur de Shapley confirme que cette diversité est utile.
\hfill\textit{→ Slides 7, 14}
\end{question}

\begin{question}
\textbf{Q : Comment choisir la valeur de $\alpha$ en pratique ?}\\[4pt]
\textbf{Réponse :} $\alpha$ est l'équivalent du paramètre $\gamma$ dans un kernel RBF classique.
En pratique : utiliser 2 valeurs par famille (petit = bande large, grand = bande étroite)
et laisser QMKL sélectionner. Les valeurs $\{0.5, 1.0, 2.5, 3.0, 4.0\}$ couvrent
la plage pertinente quand les données sont dans $[0, 2]$.
On n'optimise pas $\alpha$ par validation croisée (trop coûteux) :
la diversité des 12 kernels rend QMKL robuste au choix exact d'$\alpha$.
\hfill\textit{→ Slide 7}
\end{question}

\begin{question}
\textbf{Q : Pourquoi ne pas utiliser un kernel quantique à base de réseau neuronal (VQC) ?}\\[4pt]
\textbf{Réponse :} Les VQC (Variational Quantum Circuits) ont des paramètres entraînables,
ce qui les rend sensibles aux barren plateaus du gradient (pas seulement du kernel).
De plus, avec N=40--200 points, les gradients sont trop bruités pour converger.
VQKL (notre test variationnel) confirme cela : AUC 0.688, pire que Centered.
La PauliFeatureMap \emph{sans paramètre entraînable} est plus stable sur petits datasets.
\hfill\textit{→ Slides 8, 12}
\end{question}

% ─────────────────────────────────────────────────────────
\subsection*{C. Comparaison avec les méthodes classiques}
% ─────────────────────────────────────────────────────────

\begin{question}
\textbf{Q : Quel est l'avantage du kernel quantique par rapport à un kernel RBF classique ?}\\[4pt]
\textbf{Réponse :} Théoriquement, le kernel quantique explore un espace de Hilbert de dimension $2^Q$,
inaccessible aux kernels classiques à coût polynomial.
En pratique, sur les données testées (2026), RBF-SVM surpasse systématiquement QMKL de 5--7 pts AUC.
Il n'y a pas d'avantage mesurable aujourd'hui sur données tabulaires financières.
C'est le résultat honnête que nous présentons.
\hfill\textit{→ Slides 9, 10}
\end{question}

\begin{question}
\textbf{Q : Avez-vous testé des XGBoost ou LGBM ? QMKL serait encore plus loin.}\\[4pt]
\textbf{Réponse :} Oui, les forêts aléatoires (proches des ensembles de boosting) donnent 0.845
sur German Credit contre 0.743 pour QMKL — écart de 10 pts.
XGBoost serait probablement dans la même plage ou légèrement supérieur.
L'objectif n'est pas de battre XGBoost aujourd'hui, mais d'établir une baseline quantique rigoureuse
et d'identifier les pistes d'amélioration (QMKL-Boost, circuits adaptatifs).
\hfill\textit{→ Slides 9, 10}
\end{question}

\begin{question}
\textbf{Q : Le MKL classique (avec des kernels RBF de différentes largeurs) ne ferait-il pas mieux ?}\\[4pt]
\textbf{Réponse :} Probablement oui. Le MKL classique est une méthode mature et bien optimisée.
L'objectif de ce travail est de tester si \emph{les kernels quantiques} apportent
quelque chose de différent, pas de comparer des frameworks MKL.
Un travail futur intéressant serait de comparer le QMKL avec un MKL classique
utilisant les mêmes stratégies de combinaison (Centered Alignment sur kernels RBF).
\hfill\textit{→ Slide 20}
\end{question}

% ─────────────────────────────────────────────────────────
\subsection*{D. Stratégie QMKL : Centered Alignment et QUBO}
% ─────────────────────────────────────────────────────────

\begin{question}
\textbf{Q : Comment fonctionne concrètement le Centered Alignment ?}\\[4pt]
\textbf{Réponse :} Étant donné $M$ matrices kernel $K_1, \ldots, K_M$ et les labels $y$,
Centered Alignment maximise analytiquement la corrélation de Frobenius
$$\langle K_{\text{mkl}}, yy^\top \rangle_F = \sum_m w_m \langle K_m, yy^\top \rangle_F$$
sous contrainte $w_m \geq 0$, $\sum_m w_m = 1$.
La solution est : $w_m \propto \langle K_m, yy^\top \rangle_F$ — les poids sont proportionnels
à l'alignement de chaque kernel avec la structure de classe des données.
Aucune boucle d'optimisation, une seule multiplication matricielle.
\hfill\textit{→ Slide 8}
\end{question}

\begin{question}
\textbf{Q : Le QUBO est présenté comme quantique, mais vous le résolvez classiquement ?}\\[4pt]
\textbf{Réponse :} Oui, exactement. Avec $M=12$ kernels, l'espace de solutions est $2^{12} = 4096$ —
entièrement énumérable en quelques millisecondes sur CPU.
L'intérêt n'est pas le calcul en 2026, mais la \emph{formulation} :
QUBO est le format natif des annealers quantiques (D-Wave, Pasqal).
Quand ces machines seront suffisamment stables (horizon 2028--2030),
le même problème avec $M=10000$ kernels sera résolu en quelques microsecondes sur QPU.
\hfill\textit{→ Slides 15, 16}
\end{question}

\begin{question}
\textbf{Q : Pourquoi QUBO sélectionne-t-il seulement 2 kernels sur 12 ?}\\[4pt]
\textbf{Réponse :} C'est le résultat de l'optimisation avec $\lambda=1$ :
la pénalité de redondance est suffisamment forte pour éliminer tous les kernels corrélés
avec les 2 meilleurs (ZZ $\alpha$=2.5 et XZ $\alpha$=0.5).
Ces deux kernels sont structurellement différents : l'un capture les interactions croisées (ZZ),
l'autre les rotations mixtes (XZ). Leur complémentarité explique le gain de +4.1 pts
par rapport à la moyenne uniforme des 12.
\hfill\textit{→ Slides 15, 16}
\end{question}

\begin{question}
\textbf{Q : L'analyse Shapley sur 4095 coalitions, combien de temps ça prend ?}\\[4pt]
\textbf{Réponse :} Avec $M=12$ kernels, il y a $2^{12}-1 = 4095$ coalitions.
Pour chaque coalition, on doit calculer une matrice kernel et entraîner un SVM.
Sur CPU, cela représente environ 2--4 heures de calcul.
Avec $M=20$ kernels, ce serait $2^{20} \approx 10^6$ coalitions — prohibitif.
Pour de grands $M$, des approximations de Shapley (Monte Carlo) sont nécessaires.
\hfill\textit{→ Slide 14}
\end{question}

% ─────────────────────────────────────────────────────────
\subsection*{E. Prétraitement et pipeline de données}
% ─────────────────────────────────────────────────────────

\begin{question}
\textbf{Q : La PCA ne détruit-elle pas l'information utile pour le crédit ?}\\[4pt]
\textbf{Réponse :} Elle compresse l'information mais ne la détruit pas complètement.
Avec Q=4 composantes sur 47, on conserve $\approx 75\,\%$ de la variance totale
de German Credit. La perte de 25\,\% est une limite intrinsèque du circuit quantique :
on ne peut encoder que $Q$ valeurs. La question est donc : est-ce que le circuit
tire parti de ces 4 composantes mieux qu'un kernel RBF sur les 47 originales ?
La réponse est non, en 2026. Mais si la PCA était remplacée par un encodage adaptatif
appris (par exemple un autoencodeur), les résultats pourraient changer.
\hfill\textit{→ Slide 5}
\end{question}

\begin{question}
\textbf{Q : Pourquoi scaler vers $[0, 2]$ et pas $[0, \pi]$ ou $[-\pi, \pi]$ ?}\\[4pt]
\textbf{Réponse :} Le choix $[0, 2]$ est lié au paramètre $\alpha$ utilisé.
L'angle de rotation est $\theta = 2\alpha x_i$, donc avec $\alpha_{\max}=4.0$ et $x_i \leq 2$,
on obtient $\theta \leq 16$ — soit $\approx 2.5$ tours complets.
Ce n'est pas un choix universel : d'autres papiers utilisent $[0, \pi]$ (une demi-rotation).
L'important est que les données couvrent une plage suffisante pour que le kernel
soit informatif (pas trop concentré). Le choix $[0, 2]$ est cohérent avec la littérature
IBM/HSBC 2023 que nous reproduisons.
\hfill\textit{→ Slide 5}
\end{question}

% ─────────────────────────────────────────────────────────
\subsection*{F. Hardware IBM Torino et bruit QPU}
% ─────────────────────────────────────────────────────────

\begin{question}
\textbf{Q : Quel est le taux d'erreur typique du processeur Heron r2 ?}\\[4pt]
\textbf{Réponse :} IBM publie régulièrement les métriques de performance.
En 2025--2026, les processeurs Heron r2 atteignent :
un taux d'erreur de porte à 2 qubits (CNOT) de $\approx 0.1$--0.3\,\%,
et un taux d'erreur de lecture (readout) de $\approx 1$--3\,\%.
Pour un circuit ZZ à Q=4 avec 8 portes CNOT et 1 lecture par qubit,
l'erreur accumulée est de l'ordre de quelques pourcents — non négligeable,
surtout pour les kernels à haute concentration (ZZ $\alpha$=4).
\hfill\textit{→ Slide 18}
\end{question}

\begin{question}
\textbf{Q : Pourquoi le bruit QPU améliore-t-il certains kernels ?}\\[4pt]
\textbf{Réponse :} C'est un effet de régularisation implicite :
le bruit "lisisse" la matrice kernel en ajoutant une perturbation aléatoire.
Pour des kernels sur-concentrés (valeurs trop proches de 0), le bruit
peut paradoxalement augmenter la discrimination apparente.
Cet effet n'est cependant pas contrôlable ni reproductible :
la quantité de bruit dépend de l'état du QPU (température, calibration du jour),
ce qui rend ce "gain" inutilisable en production.
\hfill\textit{→ Slide 18}
\end{question}

\begin{question}
\textbf{Q : 1024 shots, est-ce suffisant ?}\\[4pt]
\textbf{Réponse :} Pour estimer $K(x, x') = |\langle\psi(x')|\psi(x)\rangle|^2$,
1024 shots donnent une erreur statistique de $\approx 1/\sqrt{1024} \approx 3\,\%$.
Avec des matrices kernel $30\times30$ (N=30 clients), cela représente
$30^2/2 = 450$ paires, soit $450 \times 1024 \approx 460\,000$ circuits exécutés.
Augmenter les shots à 4096 réduirait l'erreur à 1.5\,\% mais quadruplerait le coût QPU.
En pratique, 1024 shots est un compromis standard pour ce type d'expérience.
\hfill\textit{→ Slide 18}
\end{question}

% ─────────────────────────────────────────────────────────
\subsection*{G. Questions méthodologiques générales}
% ─────────────────────────────────────────────────────────

\begin{question}
\textbf{Q : 20 runs avec splits aléatoires, est-ce suffisant pour être statistiquement robuste ?}\\[4pt]
\textbf{Réponse :} Sur German Credit avec N=1000, 20 splits de 40/960 donnent
une estimation de l'AUC avec un écart-type de $\approx 0.02$--0.03.
Un test de Wilcoxon entre Centered Alignment (0.743) et RBF-SVM (0.800)
avec ces 20 mesures donne $p < 0.01$ — la différence est statistiquement significative.
20 répétitions est standard dans la littérature de comparaison de kernels quantiques
(voir IBM/HSBC 2023, MLPRAE 2025).
\hfill\textit{→ Slide 4}
\end{question}

\begin{question}
\textbf{Q : Avez-vous testé une correction d'erreur quantique (QEC) ?}\\[4pt]
\textbf{Réponse :} Non. La QEC opérationnelle nécessite des centaines de qubits physiques
par qubit logique. IBM Torino (133 qubits) n'est pas en mode QEC pour des circuits
de recherche. Nous avons utilisé les qubits physiques bruts avec atténuation d'erreurs
de base (zero-noise extrapolation non appliquée). La QEC mûre est attendue pour 2028--2030.
\hfill\textit{→ Slides 18, 20}
\end{question}

\begin{question}
\textbf{Q : Les résultats sont-ils sensibles au split train/test ?}\\[4pt]
\textbf{Réponse :} Modérément. German Credit est un dataset difficile (30\,\% de défauts,
variables hétérogènes). Avec N=40, un split particulièrement défavorable peut faire varier
l'AUC de $\pm 0.05$. C'est pourquoi les 20 répétitions sont essentielles.
Les conclusions qualitatives (Centered \textgreater{} Average, classiques \textgreater{} QMKL)
sont stables sur tous les runs.
\hfill\textit{→ Slides 9--12}
\end{question}

\begin{question}
\textbf{Q : Pourquoi utiliser N=40 points d'entraînement seulement ?}\\[4pt]
\textbf{Réponse :} Deux raisons pratiques :
\begin{enumerate}
  \item \textbf{Coût du kernel quantique :} calculer une matrice kernel $N \times N$
        nécessite $N^2/2$ évaluations de circuit. Avec N=40 : 800 circuits.
        Avec N=200 : 20\,000 circuits — 25$\times$ plus long sur simulateur, prohibitif sur QPU.
  \item \textbf{Reproductibilité hardware :} sur IBM Torino, N=30 était déjà à la limite
        du budget de temps-machine disponible pour ce projet.
\end{enumerate}
Slide 11 montre que l'écart QMKL/classique est constant de N=40 à N=200,
donc la limite N=40 ne biaise pas la conclusion.
\hfill\textit{→ Slides 4, 11}
\end{question}

\vspace{24pt}
\hrule
\vspace{8pt}
\begin{center}
  \small\color{gray}
  Document généré pour accompagner la présentation QMKL-Finance (21 slides) · Mars 2026
\end{center}

\end{document}
